\documentclass[11pt]{article}

\usepackage[margin=1in]{geometry}
\usepackage{amsmath,amssymb,amsthm,mathtools,booktabs,enumitem,microtype}
\usepackage{xcolor}
\usepackage[colorlinks=true,linkcolor=blue!50!black,urlcolor=blue!50!black,citecolor=blue!50!black]{hyperref}

\theoremstyle{plain}
\newtheorem{theorem}{Theorem}
\newtheorem{lemma}[theorem]{Lemma}
\newtheorem{proposition}[theorem]{Proposition}
\newtheorem{corollary}[theorem]{Corollary}

\theoremstyle{definition}
\newtheorem{definition}[theorem]{Definition}

\theoremstyle{remark}
\newtheorem{remark}{Remark}

\DeclareMathOperator{\conv}{conv}
\DeclareMathOperator{\sign}{sign}
\DeclareMathOperator{\dist}{dist}
\DeclareMathOperator{\rank}{rank}
\DeclareMathOperator{\Hom}{Hom}
\DeclareMathOperator{\Ext}{Ext}

\title{\bfseries Complete Proof of Beshkov's Tight Topological Expressivity Conjecture}
\author{}
\date{\vspace{-1.6em}%
\small
\begin{tabular}{@{}c@{}}
Target paper: Kosio Beshkov, \emph{Constructive Bounds for the Topological Expressivity of Neural Networks},\\
\href{https://arxiv.org/abs/2610.13931}{arXiv:2610.13931} (October 2026), Section ``Conjectured tight and lower bounds''\\[0.35em]
Lean~4 companion formalization: \texttt{lean\_proofs/Conjectures/ReluTopology.lean}\\[0.35em]
November 2026
\end{tabular}}

\begin{document}

\maketitle

\begin{abstract}
We settle Beshkov's tight topological expressivity conjecture (\href{https://arxiv.org/abs/2610.13931}{arXiv:2610.13931}) in the affirmative for all dimensions $d \ge 1$.
For the explicit one-hidden-layer ReLU network $P : \mathbb{R}^d \to \mathbb{R}^{d+1}$ of architecture $(d, d+1, d+1)$ proposed by Beshkov, we prove that $P(\mathbb{R}^d) = P(\overline{\Omega}) = \partial\Delta_0^{d+1} \cong S^d$ via an exact homeomorphism, where $\overline{\Omega} \subset [-1,2]^d$ is a compact convex polytope outside which $P$ vanishes identically; this immediately tightens the corresponding wedge-sum and product-of-spheres architectural bounds.
Conversely, we prove that no map $\Phi : X \to \mathbb{R}^{d+1}$ factoring through an affine map $\mathbb{R}^d \to \mathbb{R}^{d+1}$ (with or without output ReLU)---in particular, no $(d, d, d+1)$ ReLU network---can have image homeomorphic or homotopy equivalent to $S^d$; in fact, no continuous maps $S^d \to \Phi(X) \to S^d$ can compose to nonzero degree.
Moreover, \v{C}ech cohomology satisfies $\check{H}^d(\Phi(X); \mathbb{Z}) = 0$ for every compact topological space $X$ and continuous $h$, and singular homology satisfies $H_d(\Phi(X); G) = 0$ for every abelian group $G$ whenever $X \subseteq \mathbb{R}^d$ is a compact polyhedron, compact semialgebraic set, closed polyhedron, or $\mathbb{R}^d$ and $h$ is a ReLU network (or any continuous piecewise-linear or semialgebraic map); singular $\beta_d$ for arbitrary non-triangulable wild compacta is not covered.
\end{abstract}

\section{Exact Statement and Typographical Note}

In Section ``Conjectured tight and lower bounds'' of \href{https://arxiv.org/abs/2610.13931}{arXiv:2610.13931}, Beshkov considers the one-hidden-layer ReLU neural network $P = (p_1, \dots, p_{d+1}) : \mathbb{R}^d \to \mathbb{R}^{d+1}$ of architecture $(d, d+1, d+1)$ defined for $x = (x_1, \dots, x_d) \in \mathbb{R}^d$ by
\begin{align*}
h_i(x) &= \max(0, x_i) \qquad (i = 1, \dots, d), \\
h_{d+1}(x) &= \max\biggl(0,\, 1 - \sum_{i=1}^d x_i\biggr), \\
p_i(x) &= \max\Biggl(0,\, h_i(x) + 2\biggl(1 - \sum_{j=1}^{d+1} h_j(x)\biggr)\Biggr) \qquad (i = 1, \dots, d+1),
\end{align*}
where we correct two minor display typos in the preprint ($\sum_{i=1}^{d+1} x_i$ to $\sum_{i=1}^d x_i$ in $h_{d+1}$, and a missing closing parenthesis in $p_i$).
Beshkov conjectures that for every $d \ge 1$, the network $P$ expresses a space homeomorphic (or homotopy equivalent) to $S^d$, and that no ReLU architecture of shape $(d, d, d+1)$ can express such a sphere.

\section{Part A: Exact Homeomorphism for $(d, d+1, d+1)$}

\begin{theorem}[Part~A --- Exact Homeomorphism for $(d, d+1, d+1)$]\label{thm:partA}
For every integer $d \ge 1$, let
\[
\Delta_0^{d+1} := \conv\{0, e_1, \dots, e_{d+1}\} = \biggl\{z \in \mathbb{R}_{\ge 0}^{d+1} : \sum_{j=1}^{d+1} z_j \le 1\biggr\}
\]
be the standard $(d+1)$-simplex at the origin. Define $S(x) := \sum_{j=1}^{d+1} h_j(x)$ and $\Omega := \{x \in \mathbb{R}^d : S(x) < 2\}$. Then:
\begin{enumerate}[label=\textup{(\arabic*)}]
\item $P(\mathbb{R}^d) = \partial \Delta_0^{d+1} \cong S^d$ (exact homeomorphism).
\item $\overline{\Omega} = \{x \in \mathbb{R}^d : S(x) \le 2\}$ is a compact convex polytope satisfying
\[
\{x \in \mathbb{R}^d : P(x) \neq 0\} \subseteq \Omega \subset (-1, 2)^d,
\]
with $P \equiv 0$ outside $\Omega$. Consequently, for every domain $X$ with $\overline{\Omega} \subseteq X \subseteq \mathbb{R}^d$ (in particular $X = \overline{\Omega}$, $X = [-1,2]^d$, or $X = \mathbb{R}^d$),
\[
P(X) = P(\overline{\Omega}) = P(\mathbb{R}^d) = \partial \Delta_0^{d+1} \cong S^d.
\]
\item Therefore, the wedge-sum architectural bound tightens to
\[
\biggl(d,\, \sum_{k=1}^d \beta_k(k+1),\, \sum_{k=1}^d \beta_k(k+1)\biggr)
\]
and the product architectural bound tightens to
\[
\biggl(\sum_{n=1}^N d_n,\, \sum_{n=1}^N (d_n+1),\, \sum_{n=1}^N (d_n+1)\biggr),
\]
both realizing exact homeomorphisms to the corresponding wedge sums and products of spheres.
\end{enumerate}
\end{theorem}

\begin{proof}
Factor $P = g \circ h$, where $h : \mathbb{R}^d \to \mathbb{R}_{\ge 0}^{d+1}$ is the hidden layer and
\[
g_i(y) = \max\biggl(0,\, y_i + 2\biggl(1 - \sum_{j=1}^{d+1} y_j\biggr)\biggr) \qquad (i = 1, \dots, d+1).
\]

\paragraph{Step 1 (Hidden-layer image).}
Define
\[
H := \biggl\{y \in \mathbb{R}_{\ge 0}^{d+1} : \sum_{j=1}^{d+1} y_j \ge 1,\ \text{and}\ \biggl(\sum_{j=1}^{d+1} y_j > 1 \implies \min_{1 \le j \le d+1} y_j = 0\biggr)\biggr\}.
\]
We show that $h(\mathbb{R}^d) = H$, and that every $y \in H$ with $\sum_{j=1}^{d+1} y_j \le 2$ has a preimage $x \in \overline{\Omega}$.

\emph{Inclusion $h(\mathbb{R}^d) \subseteq H$.}
For any $x \in \mathbb{R}^d$, $y = h(x) \in \mathbb{R}_{\ge 0}^{d+1}$. Since $\max(0, u) \ge u$,
\[
\sum_{j=1}^{d+1} y_j = \sum_{i=1}^d \max(0, x_i) + \max\biggl(0,\, 1 - \sum_{i=1}^d x_i\biggr) \ge \sum_{i=1}^d x_i + \biggl(1 - \sum_{i=1}^d x_i\biggr) = 1.
\]
If $\min_j y_j > 0$, then every ReLU argument is strictly positive, so every $\max(0, \cdot)$ equals its second argument and $\sum_{j=1}^{d+1} y_j = 1$. Taking the contrapositive yields $\sum_{j=1}^{d+1} y_j > 1 \implies \min_j y_j = 0$.

\emph{Inclusion $H \subseteq h(\mathbb{R}^d)$.}
Given $y \in H$ with $S := \sum_{j=1}^{d+1} y_j \ge 1$, consider two cases:
\begin{itemize}[leftmargin=2em]
\item \emph{Case 1 ($S = 1$ or $y_{d+1} = 0$):} Set $x := (y_1, \dots, y_d) \in \mathbb{R}_{\ge 0}^d$. Then $h_i(x) = y_i$ for $1 \le i \le d$, and
\[
1 - \sum_{i=1}^d x_i = 1 - (S - y_{d+1}) = y_{d+1} - (S - 1).
\]
If $S = 1$, this equals $y_{d+1} \ge 0$; if $y_{d+1} = 0$ and $S > 1$, this equals $1 - S < 0$. In either case $h_{d+1}(x) = y_{d+1}$, so $h(x) = y$.
\item \emph{Case 2 ($S > 1$ and $y_{d+1} > 0$):} Since $y \in H$, there exists $k \in \{1, \dots, d\}$ with $y_k = 0$. Set $x_i := y_i \ge 0$ for $i \in \{1, \dots, d\} \setminus \{k\}$ and $x_k := 1 - S < 0$. Then $h_i(x) = y_i$ for $i \neq k$, $h_k(x) = 0 = y_k$, and
\[
1 - \sum_{i=1}^d x_i = 1 - (S - y_{d+1}) - (1 - S) = y_{d+1} > 0,
\]
so $h_{d+1}(x) = y_{d+1}$, giving $h(x) = y$.
\end{itemize}
Whenever $S \le 2$, this preimage satisfies $S(x) = S \le 2$, so $x \in \overline{\Omega}$.

\paragraph{Step 2 (Output-layer image).}
We show that $g(H) = g(H \cap \{S \le 2\}) = \partial \Delta_0^{d+1}$.

\emph{Inclusion $g(H) \subseteq \partial \Delta_0^{d+1}$.}
Let $y \in H$, $S = \sum_{j=1}^{d+1} y_j \ge 1$, and $t := 2(S - 1) \ge 0$, so $S = 1 + t/2$ and $z_i := g_i(y) = \max(0, y_i - t) \ge 0$.
\begin{itemize}[leftmargin=2em]
\item If $S = 1$, then $t = 0$ and $z = y \in \Delta_{\mathrm{std}}^d = \{u \in \mathbb{R}_{\ge 0}^{d+1} : \sum_{j=1}^{d+1} u_j = 1\} \subset \partial \Delta_0^{d+1}$.
\item If $S > 1$, then $t > 0$ and $\min_j y_j = 0$, hence $\min_j z_j = \max(0, -t) = 0$. Let $A := \{j : z_j > 0\} = \{j : y_j > t\}$ and $m := |A|$. If $m = 0$, then $z = 0 \in \partial \Delta_0^{d+1}$. If $m \ge 1$, then
\[
\sum_{j=1}^{d+1} z_j = \sum_{j \in A} (y_j - t) \le S - m t = 1 + \frac{t}{2} - m t = 1 - \Bigl(m - \frac{1}{2}\Bigr) t < 1,
\]
since $m - \frac{1}{2} \ge \frac{1}{2} > 0$ and $t > 0$. Because $z \in \mathbb{R}_{\ge 0}^{d+1}$, $\min_j z_j = 0$, and $\sum_{j=1}^{d+1} z_j < 1$, we have $z \in \partial \Delta_0^{d+1}$.
\end{itemize}

\emph{Inclusion $\partial \Delta_0^{d+1} \subseteq g(H \cap \{S \le 2\})$.}
Let $z \in \partial \Delta_0^{d+1}$ and $\Sigma := \sum_{j=1}^{d+1} z_j \in [0, 1]$.
\begin{itemize}[leftmargin=2em]
\item If $\Sigma = 1$: take $y := z \in H$. Then $S = 1 \le 2$, $t = 0$, and $g(y) = z$.
\item If $\Sigma = 0$ (so $z = 0$): take $y := 2 e_1$. Since $d + 1 \ge 2$, we have $y_{d+1} = 0$, $S = 2$, $y \in H$, $t = 2$, and $g(y) = 0$.
\item If $0 < \Sigma < 1$: since $z \in \partial \Delta_0^{d+1}$ and $\Sigma < 1$, we must have $\min_j z_j = 0$. Let $A := \{j : z_j > 0\}$ and $m := |A| \in \{1, \dots, d\}$. Set
\[
t := \frac{2(1 - \Sigma)}{2m - 1} > 0, \qquad
y_j := \begin{cases} z_j + t & \text{if } j \in A, \\ 0 & \text{if } j \notin A. \end{cases}
\]
Because $m \le d$, $\{1, \dots, d+1\} \setminus A \neq \varnothing$, so $\min_j y_j = 0$. Moreover,
\[
S = \sum_{j=1}^{d+1} y_j = \Sigma + m t = \Sigma + \frac{2m(1 - \Sigma)}{2m - 1} = 1 + \frac{1 - \Sigma}{2m - 1} = 1 + \frac{t}{2} \in (1, 2).
\]
Hence $y \in H \cap \{S \le 2\}$, $2(S - 1) = t$, and $g_j(y) = \max(0, y_j - t) = z_j$ for all $j = 1, \dots, d+1$.
\end{itemize}

\paragraph{Step 3 (Compact convex support).}
The function $S(x) = \sum_{i=1}^d x_i^+ + \bigl(1 - \sum_{i=1}^d x_i\bigr)^+$ is convex and continuous with $S(0) = 1 < 2$, so $\Omega = \{S < 2\}$ is nonempty and convex and its closure is $\overline{\Omega} = \{S \le 2\}$, which is closed and convex. For every $k \in \{1, \dots, d\}$:
\begin{itemize}[leftmargin=2em]
\item $S(x) \ge x_k^+ \ge x_k$, so $S(x) < 2 \implies x_k < 2$;
\item writing $x_i = x_i^+ - x_i^-$, we have
\[
S(x) \ge \sum_{i=1}^d x_i^+ + \biggl(1 - \sum_{i=1}^d x_i^+ + \sum_{i=1}^d x_i^-\biggr) = 1 + \sum_{i=1}^d x_i^- \ge 1 + x_k^-,
\]
so $S(x) < 2 \implies x_k > -1$.
\end{itemize}
Hence $\Omega \subset (-1, 2)^d$ and $\overline{\Omega} \subset [-1, 2]^d$ is a compact convex polytope. Whenever $S(x) \ge 2$, we have
\[
t = 2(S(x) - 1) = S(x) + (S(x) - 2) \ge S(x) \ge h_i(x)
\]
for every $i$, so $P(x) = 0$. Combining Steps~1--3 gives $P(X) = P(\overline{\Omega}) = P(\mathbb{R}^d) = \partial \Delta_0^{d+1} \cong S^d$ for every $\overline{\Omega} \subseteq X \subseteq \mathbb{R}^d$.

\paragraph{Step 4 (Wedge sums and products).}
For wedge sums, the shifted networks $x \mapsto P^{(k)}(x_1 - 3r, x_2, \dots, x_k)$ ($r = 0, 1, \dots, N-1$) have pairwise disjoint open supports contained in $(3r - 1, 3r + 2) \times (-1, 2)^{d-1}$ and vanish identically outside them. Placing their outputs in disjoint coordinate blocks glues their images $\partial \Delta_0^{k+1} \cong S^k$ exclusively at the origin $0$, yielding an exact homeomorphism to $\bigvee_{k=1}^d \bigl(\bigvee^{\beta_k} S^k\bigr)$ with architecture $\bigl(d, \sum_{k=1}^d \beta_k(k+1), \sum_{k=1}^d \beta_k(k+1)\bigr)$.
Stacking $N$ independent copies on disjoint input blocks yields $\prod_{n=1}^N \partial \Delta_0^{d_n+1} \cong \prod_{n=1}^N S^{d_n}$ with architecture $\bigl(\sum_{n=1}^N d_n, \sum_{n=1}^N (d_n+1), \sum_{n=1}^N (d_n+1)\bigr)$.
\end{proof}

\section{Part B: Impossibility for $(d, d, d+1)$}

\begin{theorem}[Part~B --- Impossibility for $(d, d, d+1)$]\label{thm:partB}
Let $d \ge 1$, let $X$ be any topological space, and let $\Phi : X \to \mathbb{R}^{d+1}$ be any map of the form
\[
\Phi(x) = R(W_2 h(x) + b_2) \qquad \text{or} \qquad \Phi(x) = W_2 h(x) + b_2,
\]
where $R(u) = \max(0, u)$ is componentwise ReLU, $W_2 \in \mathbb{R}^{(d+1)\times d}$, $b_2 \in \mathbb{R}^{d+1}$, and $h : X \to \mathbb{R}^d$ is arbitrary (in particular, any $(d, d, d+1)$ ReLU network, or any deeper network whose final hidden layer has width at most $d$). Then:
\begin{enumerate}[label=\textup{(\arabic*)}]
\item $\Phi(X)$ is neither homeomorphic nor homotopy equivalent to $S^d$. More strongly, there do not exist continuous maps $f : S^d \to \Phi(X)$ and $g : \Phi(X) \to S^d$ with $\deg(g \circ f) \neq 0$.
\item For every compact topological space $X$ and continuous $h$, the \v{C}ech cohomology vanishes: $\check{H}^d(\Phi(X); \mathbb{Z}) = 0$. Moreover, whenever $X \subseteq \mathbb{R}^d$ is a compact polyhedron or compact semialgebraic set, or $X = \mathbb{R}^d$ (or any closed polyhedron in $\mathbb{R}^d$), and $h$ is a ReLU network (or any continuous piecewise-linear or semialgebraic map), singular homology also vanishes:
\[
H_d(\Phi(X); G) = 0 \qquad \text{for every abelian coefficient group } G.
\]
(For arbitrary wild non-triangulable compact domains $X$, singular $H_d(\Phi(X))$ is not controlled by Alexander duality alone.)
\end{enumerate}
\end{theorem}

Since $W_2 \in \mathbb{R}^{(d+1)\times d}$ has rank at most $d$, there exists $n \in \mathbb{R}^{d+1} \setminus \{0\}$ with $W_2^\top n = 0$. Setting $\alpha := n \cdot b_2$, the affine image $W_2 \mathbb{R}^d + b_2$ lies inside the affine hyperplane
\[
H := \{u \in \mathbb{R}^{d+1} : n \cdot u = \alpha\}.
\]
Thus $\Phi(X) \subseteq R(H)$ (with output ReLU) or $\Phi(X) \subseteq H$ (linear output). The key geometric ingredient is the following lemma.

\begin{lemma}[Straight-Ray Escape Lemma]\label{lem:escape}
Let $H = \{u \in \mathbb{R}^{d+1} : n \cdot u = \alpha\}$ be an affine hyperplane ($n \in \mathbb{R}^{d+1} \setminus \{0\}$, $\alpha \in \mathbb{R}$). For every $p \in \mathbb{R}^{d+1} \setminus R(H)$, there exists a straight ray $\gamma(t) = p + t v$ ($t \ge 0$, $v \in \mathbb{R}^{d+1} \setminus \{0\}$) contained in $\mathbb{R}^{d+1} \setminus R(H)$. Similarly, for every $p \in \mathbb{R}^{d+1} \setminus H$, the ray $p + t\,\sign(n \cdot p - \alpha)\,n$ ($t \ge 0$) stays in $\mathbb{R}^{d+1} \setminus H$.
\end{lemma}

\begin{proof}
The linear-hyperplane claim is immediate because $n \cdot (p + t\,\sign(n \cdot p - \alpha)\,n) - \alpha = (n \cdot p - \alpha) + t\,\sign(n \cdot p - \alpha)\,\|n\|_2^2$ never vanishes for $t \ge 0$.

Now fix $p \in \mathbb{R}^{d+1} \setminus R(H)$.
\begin{itemize}[leftmargin=2em]
\item \emph{Case 1 ($\exists k$ with $p_k \le 0$):} Take $\gamma(t) := p - t e_k$ ($t \ge 0$). At $t = 0$, $\gamma(0) = p \notin R(H)$. For every $t > 0$, $\gamma(t)_k = p_k - t < 0$, so $\gamma(t) \notin \mathbb{R}_{\ge 0}^{d+1} \supseteq R(H)$.
\item \emph{Case 2 ($p \in \mathbb{R}_{>0}^{d+1}$):} Since $R(p) = p \notin R(H)$, we have $p \notin H$, so $n \cdot p \neq \alpha$. Replacing $(n, \alpha)$ with $(-n, -\alpha)$ if necessary, assume $n \cdot p < \alpha$.
  \begin{itemize}[leftmargin=1.5em]
  \item \emph{Subcase 2a ($\exists k$ with $n_k > 0$):} Take $\gamma(t) := p - t e_k$ ($t \ge 0$).
    \begin{itemize}[leftmargin=1.2em]
    \item For $0 \le t < p_k$: $\gamma(t) \in \mathbb{R}_{>0}^{d+1}$, so its unique $R$-preimage is $u = \gamma(t)$, which satisfies $n \cdot u = n \cdot p - t n_k \le n \cdot p < \alpha$, hence $\gamma(t) \notin R(H)$.
    \item For $t = p_k$: $\gamma(p_k)_k = 0$ and $\gamma(p_k)_j = p_j > 0$ for $j \neq k$. Any $u \in \mathbb{R}^{d+1}$ with $R(u) = \gamma(p_k)$ must satisfy $u_j = p_j$ ($j \neq k$) and $u_k \le 0 < p_k$. Since $n_k > 0$,
    \[
    n \cdot u = n \cdot p + n_k (u_k - p_k) < n \cdot p < \alpha,
    \]
    so $u \notin H$; hence $\gamma(p_k) \notin R(H)$.
    \item For $t > p_k$: $\gamma(t)_k = p_k - t < 0$, so $\gamma(t) \notin \mathbb{R}_{\ge 0}^{d+1} \supseteq R(H)$.
    \end{itemize}
  \item \emph{Subcase 2b ($n_i \le 0$ for all $i$):} Since $n \neq 0$, pick $k$ with $n_k < 0$ and take $\gamma(t) := p + t e_k$ ($t \ge 0$). For all $t \ge 0$, $\gamma(t) \in \mathbb{R}_{>0}^{d+1}$, so its unique $R$-preimage is $u = \gamma(t)$ with $n \cdot u = n \cdot p + t n_k \le n \cdot p < \alpha$, hence $\gamma(t) \notin R(H)$. \qedhere
  \end{itemize}
\end{itemize}
\end{proof}

\begin{proof}[Proof of Theorem~\ref{thm:partB}]
\textbf{Step 1 (Non-separation of every compact $K \subseteq R(H)$).}
Decomposing $\mathbb{R}^{d+1} = \bigcup_{\sigma \in \{-,+\}^{d+1}} Q_\sigma$ into closed orthants, $R|_{Q_\sigma}$ is a linear coordinate projection $\pi_\sigma$, so $R(H) = \bigcup_\sigma \pi_\sigma(H \cap Q_\sigma)$ is a finite union of closed convex polyhedra, each contained in an affine subspace $\pi_\sigma(H)$ of dimension at most $d$. Hence $R(H)$ has empty interior in $\mathbb{R}^{d+1}$.

Let $K \subseteq R(H)$ be compact and pick $R_0 > 0$ with $K \subset B(0, R_0)$. Since $d + 1 \ge 2$, the exterior $E := \mathbb{R}^{d+1} \setminus \overline{B(0, R_0)}$ is path-connected. Given any $x \in \mathbb{R}^{d+1} \setminus K$, choose $\varepsilon > 0$ with $B(x, \varepsilon) \subseteq \mathbb{R}^{d+1} \setminus K$. Since $R(H)$ has empty interior, pick $p \in B(x, \varepsilon) \setminus R(H)$. Walking along the segment $[x, p] \subset B(x, \varepsilon)$ and then along the escape ray $\gamma$ of Lemma~\ref{lem:escape} from $p$ connects $x$ to $E$ in $\mathbb{R}^{d+1} \setminus K$. Thus $\mathbb{R}^{d+1} \setminus K$ is path-connected, and hence $S^{d+1} \setminus K$ is path-connected in the one-point compactification $S^{d+1} = \mathbb{R}^{d+1} \cup \{\infty\}$. The case $K \subseteq H$ is identical.

\paragraph{Step 2 (Proof of (1) via Tietze extension and Alexander duality).}
Suppose for contradiction that continuous maps $f : S^d \to \Phi(X)$ and $g : \Phi(X) \to S^d$ satisfy $\deg(g \circ f) \neq 0$. Let $Z := f(S^d) \subseteq \Phi(X) \subseteq R(H)$ (or $\subseteq H$), which is nonempty and compact. Viewing $S^d \subset \mathbb{R}^{d+1}$, extend $g|_Z : Z \to S^d$ componentwise by the Tietze Extension Theorem to a continuous map $G : \mathbb{R}^{d+1} \to \mathbb{R}^{d+1}$. Let $U := \{y \in \mathbb{R}^{d+1} : G(y) \neq 0\}$ (an open neighborhood of $Z$) and define $\tilde{g} : U \to S^d$ by $\tilde{g}(y) := G(y)/\|G(y)\|_2$, so $\tilde{g}|_Z = g|_Z$.

Choose $M > 0$ with $Z \subset (-M, M)^{d+1}$, and set $P_M := R(H) \cap [-M, M]^{d+1}$ (or $H \cap [-M, M]^{d+1}$), a finite union of compact convex polytopes of dimension at most $d$, hence triangulable as a finite simplicial complex of dimension at most $d$ (Rourke--Sanderson, \emph{Introduction to Piecewise-Linear Topology}, Theorem~2.11). Taking iterated barycentric subdivisions so every simplex has diameter strictly less than $\dist(Z, \mathbb{R}^{d+1} \setminus U)$ (or taking $N = P_M$ if $U = \mathbb{R}^{d+1}$), let $N$ be the union of all closed simplices meeting $Z$. Then $N$ is a nonempty compact finite simplicial complex (hence compact and locally contractible) with $Z \subseteq N \subseteq R(H) \cap U$ (respectively $H \cap U$).

By Step~1, $\tilde{H}_0(S^{d+1} \setminus N; \mathbb{Z}) = 0$, so by Alexander Duality (Hatcher, \emph{Algebraic Topology}, Theorem~3.44),
\[
H^d(N; \mathbb{Z}) \cong \tilde{H}_0(S^{d+1} \setminus N; \mathbb{Z}) = 0.
\]
Yet the composite $S^d \xrightarrow{f} N \xrightarrow{\tilde{g}|_N} S^d$ equals $g \circ f$, so $(g \circ f)^* : H^d(S^d; \mathbb{Z}) \to H^d(S^d; \mathbb{Z})$ is multiplication by $\deg(g \circ f) \neq 0$ while factoring through $H^d(N; \mathbb{Z}) = 0$, a contradiction.

\paragraph{Step 3 (Proof of (2) --- vanishing of $\check{H}^d$ and $H_d(\Phi(X); G)$).}
For any compact $K = \Phi(X) \subseteq R(H)$ (or $\subseteq H$) with $h$ continuous, Step~1 gives $\tilde{H}_0(S^{d+1} \setminus K; \mathbb{Z}) = 0$, so \v{C}ech Alexander duality immediately yields $\check{H}^d(K; \mathbb{Z}) \cong \tilde{H}_0(S^{d+1} \setminus K; \mathbb{Z}) = 0$.

When $N := \Phi(X) \subseteq R(H)$ (or $\subseteq H$) is a compact polyhedron (triangulable by Rourke--Sanderson, Theorem~2.11) or a compact semialgebraic set (triangulable by Hardt/\L{}ojasiewicz semialgebraic triangulation, hence compact and locally contractible of dimension at most $d$), Alexander duality gives $H^d(N; \mathbb{Z}) = 0$. By the Universal Coefficient Theorem for cohomology,
\[
\Hom_{\mathbb{Z}}(H_d(N; \mathbb{Z}), \mathbb{Z}) = 0
\qquad \text{and} \qquad
\Ext^1_{\mathbb{Z}}(H_{d-1}(N; \mathbb{Z}), \mathbb{Z}) = 0.
\]
Since $\dim N \le d$, $H_d(N; \mathbb{Z}) = \ker(\partial_d)$ is free abelian, forcing $H_d(N; \mathbb{Z}) = 0$, and $H_{d-1}(N; \mathbb{Z})$ is finitely generated and torsion-free. By the Universal Coefficient Theorem for homology, $H_d(N; G) = 0$ for every abelian group $G$.

Finally, if $X = \mathbb{R}^d$ (or any closed polyhedron in $\mathbb{R}^d$), write $\Phi(X) = \bigcup_{M=1}^\infty Z_M$ with $Z_M := \Phi(X) \cap [-M, M]^{d+1}$, each a compact polyhedron of dimension at most $d$ inside $R(H)$ (or $H$). Because singular homology commutes with directed unions of compact subsets covering all compacta,
\[
H_d(\Phi(X); G) \cong \varinjlim_M H_d(Z_M; G) = 0. \qedhere
\]
\end{proof}

\section{Structural Remarks for Arbitrary Output Width $(d, d, m)$}

We record three brief structural observations regarding networks of architecture $(d, d, m)$ with arbitrary output width $m$.

\begin{remark}[Full contractibility when $\rank(W_1) = d$ and $X = \mathbb{R}^d$; sketch / partial]
If $W_1 \in \mathbb{R}^{d \times d}$ is invertible, then $h(\mathbb{R}^d) = \mathbb{R}_{\ge 0}^d$ is convex, so $C = W_2 \mathbb{R}_{\ge 0}^d + b_2 \subset \mathbb{R}^m$ is a closed convex polyhedral cone. Whenever $C$ meets the nonpositive orthant $\mathbb{R}_{\le 0}^m$ at some $z_0 \le 0$, the straight-line homotopy $(z, t) \mapsto R_m((1-t)z + t z_0)$ contracts $R_m(C)$ to $0$. The general case $C \cap \mathbb{R}_{\le 0}^m = \varnothing$ is left conjectural.
\end{remark}

\begin{remark}[Sharpness of $\rank(W_1) = d$]
If $\rank(W_1) < d$, lower-dimensional spheres can appear. For example, for architecture $(2, 2, 3)$ with $h(x_1, x_2) = (\max(0, x_1),\, \max(0, 1 - x_1))$, applying Theorem~\ref{thm:partA} for $d = 1$ yields $S^1 \subset \mathbb{R}^3$ (so $H_1 \neq 0$ while $H_2 = 0$, confirming that Theorem~\ref{thm:partB} requires no rank hypothesis on $W_1$).
\end{remark}

\begin{remark}[Compact convex $X \subset \mathbb{R}^d$ and arbitrary output width $m$; conjectural sketch]
Splitting $R_d(W_1 X + b_1) = A \cup B$ into its compact convex interior-orthant part $A$ and boundary part $B \subseteq \partial \mathbb{R}_{\ge 0}^d$ ($\dim B \le d - 1$) suggests $\check{H}^d(\Phi(X); G) = 0$ via Vietoris--Begle and strong excision; this extension is left unproved here.
\end{remark}

\section{Formalization, Verification, and Attribution}

\paragraph{Lean~4 formalization.}
In \texttt{lean\_proofs/Conjectures/ReluTopology.lean} (compiled under Lean~\texttt{v4.24.0} with Mathlib, zero \texttt{sorry}s, and standard axioms only), we machine-check:
\begin{enumerate}[label=\textup{(\arabic*)}]
\item \emph{Theorem~\ref{thm:partA} (end-to-end):} Continuity of $P$ (\texttt{continuous\_P\_net}), exact image $P(\mathbb{R}^d) = P(\overline{\Omega}) = \partial \Delta_0^{d+1}$ (\texttt{beshkov\_part\_A}), compactness, convexity, and box containment of $\overline{\Omega} \subseteq [-1, 2]^d$ (\texttt{isCompact\_omegaBar}, \texttt{convex\_omegaBar}, \texttt{omegaBar\_\allowbreak subset\_box}), vanishing outside $\Omega$ (\texttt{P\_net\_eq\_\allowbreak zero\_of\_S\_ge\_two}), and the explicit radial homeomorphism $\partial \Delta_0^{d+1} \cong S^d$ from the barycenter $(\frac{1}{d+2}, \dots, \frac{1}{d+2})$ (\texttt{boundarySimplex\_\allowbreak homeomorph\_sphere}).
\item \emph{Theorem~\ref{thm:partB} (geometric engine and impossibility for geometric spheres and simplex boundaries):} Hyperplane containment (\texttt{bottleneck\_in\_\allowbreak hyperplane}, \texttt{bottleneck\_relu\_in\_\allowbreak relu\_hyperplane}), the Straight-Ray Escape Lemma (\texttt{straight\_ray\_escape}), and the impossibility of containing any Euclidean sphere (\texttt{bottleneck\_not\_\allowbreak superset\_sphere}) or $\partial \Delta_0^{d+1}$ (\texttt{bottleneck\_not\_\allowbreak superset\_\allowbreak boundarySimplex}).
\end{enumerate}
Extending Theorem~\ref{thm:partB} from geometric spheres and simplex boundaries to arbitrary wild topological embeddings of $S^d$ and cohomology/homology vanishing uses Alexander duality and \v{C}ech cohomology (not yet in Mathlib) and is proved in prose in Section~3 above. Companion numerical checks are provided in \texttt{solutions/relu\_topology/verify\_conjecture.py}.

\paragraph{Attribution and priority note.}
No completed proof has appeared in the literature (as of November~26, 2026, \href{https://arxiv.org/abs/2610.13931}{arXiv:2610.13931} remains v1). An independent public preregistration (\href{https://github.com/the-omega-institute/trureturing/issues/25155}{\texttt{the-omega-\allowbreak institute/\allowbreak trureturing\#25155}}, opened November~1, 2026, open with no completed proof or Lean file) outlines a similar Part~A construction, whereas the Part~B homotopy-theoretic argument via Straight-Ray Escape and Alexander Duality---ruling out not only homeomorphisms, but all degree-nonzero maps $S^d \to \Phi(X) \to S^d$ and all homotopy equivalences for arbitrary $X$ and arbitrary $h$---is completely independent and strictly stronger.

\end{document}
